% Enunciados selecionados - Capítulos 3, 4 e 5 (Português)
\section*{Enunciados selecionados}

\textbf{Capítulo 3}
\begin{enumerate}
  \item[3.3] Como a área de superfície e o volume de uma esfera são dimensionados? Por que? (Dica: analise esferas de raios 1 e $R$.)
  \item[3.6] Mostre como a equação $y = m x^b$ se torna uma equação linear em um gráfico log-log.
  \item[3.9] Use a análise dimensional para determinar que a potência é igual à taxa de variação da energia cinética.
  \item[3.18] Confirme que a eq. (3.28) está correta.
  \item[3.19] Mostre que a deflexão $w_{mp}$ de uma viga com rigidez à flexão $EI$, comprimento $L$ e sob uma carga concentrada $P$ é governada pelos dois grupos adimensionais na eq. (3.34).
  \item[3.28] Usando um raciocínio semelhante ao que nos levou à eq. (3.13), mostre que a velocidade máxima na qual os animais podem correr é independente do tamanho.
  \item[3.32] Descobriu-se que um certo pepino tinha células que se dividiam quando cresciam até 1,5 vezes o volume das células "remanescentes". As células normalmente se dividem de modo que a proporção entre sua superfície e sua massa permaneça constante. O pepino é descrito como um pepino "normal"?
  \item[3.35] Quando os elementos estruturais denominados vigas vibram... (Nota: o enunciado original está cortado/incompleto.)
\end{enumerate}

\textbf{Capítulo 4}
\begin{enumerate}
  \item[4.1] Mostre que a eq. (4.7) pode ser obtida substituindo $ix$ por $x$ na eq. (4.6).
  \item[4.2] Determine os primeiros quatro termos das expansões de Taylor de $\tan x$ e $\cot x$ sobre $x = 0$.
  \item[4.3] Determine os quatro primeiros termos das expansões de Taylor de $\tanh x$ e $\coth x$ sobre $x = 0$.
  \item[4.4] Use a análise dimensional para determinar como o parâmetro catenário, $c$, está relacionado ao componente horizontal constante da tensão do cabo, $T_0$, e seu peso por unidade de comprimento, $\gamma$.
  \item[4.6] O que pesa um corpo que pesa 10 N na superfície da Terra a uma altura de 10 m? No pico do Monte Everest? (Use dados planetários, se necessário.)
  \item[4.10] Se o potencial gravitacional for $V_g = -G M_E m / R$, encontre a expressão exata do potencial em função da altitude $z$ acima da superfície.
  \item[4.11] Faça a expansão binomial do resultado do Problema 4.10 para obter a energia potencial à primeira ordem em $z$.
  \item[4.12] Preencha a tabela para a ordem de dois termos: $\sin x$, $\cos x$, $1-\sin x$, $1-\cos x$.
  \item[4.13] Preencha a tabela para a ordem de dois termos: $\sinh x$, $\cosh x$, $1-\sinh x$, $1-\cosh x$.
  \item[4.14] Desenvolva o coeficiente de expansão volumétrica $\beta$ para um sólido de dimensões $L_0, W_0, H_0$ em paralelo ao coeficiente de superfície $\gamma$ da eq. (4.37).
  \item[4.17] Arredonde cada número para duas cifras significativas: (a) 5,237 (b) 0,82549 (c) 81,356 (d) $\pi$ (e) 6,2305 (f) 0,0428 (g) 10,45 (h) 4,035.
  \item[4.18] Arredonde cada número para três cifras significativas: (a) 5,237 ... (h) 4,035.
  \item[4.19] Complete multiplicações e expresse com algarismos significativos corretos.
  \item[4.20] 99.9 e 100.1 têm o mesmo número de algarismos significativos? Explique.
  \item[4.21] Estime os intervalos para (a) 7.7 (b) 7.70 (c) 1200 (d) $1.200\times10^{-3}$.
  \item[4.22] Em que porcentagem mudaria o período de um pêndulo se seu comprimento fosse reduzido pela metade; em um terço; a um terço do original?
  \item[4.24] Como mudaria o período de um pêndulo em Marte? Em Plutão?
  \item[4.35] Estime o erro da aproximação $\sin x$ por Taylor de ordem $n=4$ calculando o resto $R_5$.
  \item[4.36] Confirmar se $\sin x \simeq x$ e $\tan x \simeq x$ dão aproximações semelhantes.
  \item[4.37] Voltímetro analógico: corrija leituras e calcule erros percentuais para valores dados.
  \item[4.39] (a) Série de Taylor de $e^x$ em $0$. (b) Calcule $e^{0.5}$ com quatro termos para cinco algarismos significativos.
  \item[4.45] Analise dados experimentais $V$ vs $I$; discuta erros e faça um ajuste visual.
  \item[4.46] Obtenha a reta por mínimos quadrados para os mesmos dados e compare.
  \item[4.50] Aproximação para $f(r)=\sqrt{a^2+r^2}$ para $r\gg a$ (expansão binomial).
\end{enumerate}

\textbf{Capítulo 5}
\begin{enumerate}
  \item[5.8] Que taxa anual seria necessária para dobrar o dinheiro em sete anos?
  \item[5.9] Verifique que a eq. (5.12) resolve a equação diferencial exponencial dada em (5.10) usando $\int du/u = \ln u + C$.
  \item[5.35] Obtenha os números reais da população mundial para 1970, 1980, 1990 e 2000 e atualize as Figuras 5.1 e 5.2.
  \item[5.40] Quanto reservar em 2002 a 5.5% a.a. para acumular US\$1.000.000 até 2022? E até 2042?
  \item[5.47] Batalha de Iwo Jima: com $F_0=54000$, $E_0=21500$, $\lambda_F=0.0106$, $\lambda_E=0.0544$ por dia — quanto duraria e quantos sobreviveriam ao final?
\end{enumerate}

% Fim do arquivo de enunciados (Português)
\begin{enumerate}
    \item[\textbf{Problema 3.1.}] Com base em que os liliputianos concluíram que Gulliver precisava de 1728 vezes mais comida do que eles?

    \item[\textbf{Problema 3.2.}] Como a conclusão dos liliputianos mudaria se tivessem pensado na troca de energia entre uma pessoa (de qualquer tamanho) e o ambiente circundante?

    \item[\textbf{Problema 3.3.}] Como a área de superfície e o volume de uma esfera são dimensionados? Por que? (Dica: analise esferas de raios 1 e R.)

    \item[\textbf{Problema 3.4.}] Explique o que aconteceria com um ângulo entre duas linhas, inscrito num balão enquanto ele era inflado para um raio R a partir de um raio de 1.

    \item[\textbf{Problema 3.5.}] Confirme que a eq. (3.2) retrata adequadamente a linha reta traçada na Figura 3.2.

    \item[\textbf{Problema 3.6.}] Mostre como a equação $y = mx^b$ se torna uma equação linear em um gráfico log-log.

    \item[\textbf{Problema 3.7.}] Escreva a eq. (3.6) em uma forma adequada para um gráfico log-log.

    \item[\textbf{Problema 3.8.}] Defina a relação dimensional da eq. (3.10).

    \item[\textbf{Problema 3.9.}] Use a análise dimensional para determinar que a potência é igual à taxa de variação da energia cinética.

    \item[\textbf{Problema 3.10.}] Defina a relação dimensional da eq. (3.12).

    \item[\textbf{Problema 3.11.}] Confirme que a eq. (3.13) mostra que $v$ é independente de $L$.

    \item[\textbf{Problema 3.12.}] Qual é o alcance auditivo de um elefante? Uma baleia? Como esses intervalos se comparam aos dos humanos?

    \item[\textbf{Problema 3.13.}] Considere que as dimensões da eq. (3.15) são 1/T.

    \item[\textbf{Problema 3.14.}] Confirme que as dimensões da eq. (3.16) são 1/T.

    \item[\textbf{Problema 3.15.}] Entende-se que as dimensões da eq. (3.17) são 1/T.

    \item[\textbf{Problema 3.16.}] Sob quais condições a eq. (3.24) é dimensionalmente consistente?

    \item[\textbf{Problema 3.17.}] Confirme que a tensão da eq. (3.27) satisfaz a eq. (3.24).

    \item[\textbf{Problema 3.18.}] Confirme que a eq. (3.28) está correta.

    \item[\textbf{Problema 3.19.}] Mostre que a deflexão $w_{mp}$ de uma viga com rigidez à flexão $EI$, comprimento $L$ e sob uma carga concentrada $P$ é governada pelos dois grupos adimensionais na eq. (3.34).

    \item[\textbf{Problema 3.20.}] Por que o torque, $\tau$, no aparelho da Figura 3.12 é uma constante?

    \item[\textbf{Problema 3.21.}] Expresso em termos das propriedades geométricas e gravitacionais da roda, qual é a magnitude do torque no Problema 3.20?

    \item[\textbf{Problema 3.22.}] Confirme que a eq. (3.30) é dimensionalmente correta.

    \item[\textbf{Problema 3.23.}] Confirme que a eq. (3.31) é dimensionalmente correta.

    \item[\textbf{Problema 3.24.}] Calcule e confirme o número estimado de rotações na última coluna do quadro 3.1.

    \item[\textbf{Problema 3.25.}] Confirme que a eq. (3.39) é a representação correta da eq. (3.37) em termos dos cinco fatores de escala de uma viga simples.

    \item[\textbf{Problema 3.26.}] Formule uma hipótese para explicar por que um pombo-torcaz e um urubu parecem ter proporções tão diferentes de $W_{fm} / W_b$ na Figura 3.2.

    \item[\textbf{Problema 3.27.}] Mostre que a equação que descreve o gráfico log-log da Figura 3.7 pode ser encontrada como $h \simeq 1{,}23l^{0{,}68}$, onde $h$ e $l$ são, respectivamente, a altura da nave e o comprimento da igreja tornados adimensionais dividindo cada um por 1 pé.

    \item[\textbf{Problema 3.28.}] Usando um raciocínio semelhante ao que nos levou à eq. (3.13), mostre que a velocidade máxima na qual os animais podem correr é independente do tamanho.

    \item[\textbf{Problema 3.29.}] A velocidade do sangue na aorta está relacionada à diferença de pressão entre o coração e as artérias. Encontre a relação entre a velocidade do sangue e a diferença de pressão. (Dica: use o teorema da energia do trabalho como fizemos para o pássaro pairando na Seção 3.3.2.)

    \item[\textbf{Problema 3.30.}] O pernilongo, um pequeno pássaro de pernas compridas, foi descrito em \textit{As Viagens de Gulliver} como pesando 4,5 onças e tendo pernas de 8 polegadas de comprimento. Um flamingo tem uma forma semelhante e pesa 4 libras. Aplique argumentos de escala para mostrar que as pernas devem ter cerca de 20 polegadas de comprimento (como realmente são!).

    \item[\textbf{Problema 3.31.}] Dado que um robin pesa cerca de 2 onças, poderíamos dimensionar o comprimento de suas pernas a partir dos dados de pernas de pau fornecidos no Problema 3.30? Explique sua resposta.

    \item[\textbf{Problema 3.32.}] Descobriu-se que um certo pepino tinha células que se dividiam quando cresciam até 1,5 vezes o volume das células remanescentes. As células normalmente se dividem de modo que a proporção entre sua superfície e sua massa permaneça constante. O pepino é descrito como um pepino normal?

    \item[\textbf{Problema 3.33.}] Encontrar o intervalo de valores da variável $x$ para o qual a aproximação $\cosh(x/\lambda) \simeq \frac{1}{2}e^{x/\lambda}$ é aceitável, para fatores de escala $\lambda = 1$ e $\lambda = 6$. Plote ambas as funções para cada um dos dois fatores de escala.

    \item[\textbf{Problema 3.34.}] Um experimento para determinar o período natural ou fundamental de oscilação de um sistema simples de massa-mola (ver figura P3.34) é estabelecido da seguinte forma. Uma mola de rigidez $k$ é fixada em uma extremidade e conectada a uma massa $m$ na outra, com a massa sendo capaz de se mover ao longo de uma pista de ar ideal e sem atrito. A massa é deslocada a uma distância $x_0$ de sua posição inicial de repouso, após o que oscila ao longo da pista de ar em torno dessa posição inicial. O tempo necessário para uma oscilação completa, isto é, o período, é medido várias vezes por vários períodos sucessivos, com os resultados sendo comparados com a fórmula teórica para o período, $T$: $T = 2\pi\sqrt{m/k}$. Supondo que $k$ e $m$ sejam conhecidos e que o cronômetro usado para medir o período tenha precisão de $\pm 1\%$. Quais são as possíveis armadilhas que podem impedir a determinação experimental bem-sucedida de $T$?

    \item[\textbf{Problema 3.35.}] Quando os elementos estruturais denominados vigas vibram... \textit{(Nota: Este problema está cortado no livro original, faltando o restante da frase que indicaria como as vibrações dependem das propriedades da viga).}

    \item[\textbf{Problema 3.36.}] Uma viga de aço de 20 cm de comprimento deve ser usada para modelar um protótipo de viga de madeira cujo vão é de 3,6 m. (a) Verifique se o grupo adimensional que contém a carga, o módulo e o comprimento é $P/(EL^2)$. (b) Se a viga de madeira deve suportar uma carga de 9000 N em um ponto a 1,5 m da extremidade esquerda, que carga deve ser aplicada ao modelo para determinar se o protótipo pode suportar a carga pretendida? (Suponha que a capacidade de carga seja o único comportamento de interesse aqui.)

\end{enumerate}
